\chapter{Charakterisierung, Existenz und Gleichheit}
\label{b05guide:ch:criteria}

\section{Was charakterisiert eine Funktion?}

\begin{theorem}[Eindeutige Zuordnung und einelementige Wertemengen]
\label{b05guide:thm:criteria}
Sei $R\subseteq A\times B$. Die folgenden Bedingungen sind gleichwertig:
\begin{enumerate}
 \item $R\colon A\to B$ ist eine Funktion.
 \item Für jedes $x\in A$ existiert genau ein $y\in B$ mit $(x,y)\in R$.
 \item Für jedes $x\in A$ ist
       $\{y\in B\mid(x,y)\in R\}$ eine einelementige Menge.
\end{enumerate}
\end{theorem}

\begin{theorem}[Funktionen als minimale totale Relationen]
\label{b05guide:thm:minimal}
Sei $R$ eine totale Relation von $A$ nach $B$ im Sinn von Band~4.
Dann gilt
\[
 \begin{gathered}
 R\colon A\to B\\
 \Longleftrightarrow\quad
 \forall S\subseteq R\,
 \bigl(S\text{ ist total von }A\text{ nach }B\Rightarrow S=R\bigr).
 \end{gathered}
\]
Eine totale Relation ist also genau dann eine Funktion,
wenn keine echte Teilmenge ihres Graphen auf demselben
Definitionsbereich total bleibt.
\end{theorem}

\begin{theorem}[Charakterisierung durch Urbilder von Komplementen]
\label{b05guide:thm:complementcriterion}
Sei $R\subseteq A\times B$. Für $V\subseteq B$ setze
\[
 R^{-1}[V]=\{x\in A\mid\exists y\in V\,((x,y)\in R)\}.
\]
Dann gilt
\[
 R\colon A\to B
 \quad\Longleftrightarrow\quad
 \forall V\subseteq B\,
   \bigl(R^{-1}[B\setminus V]=A\setminus R^{-1}[V]\bigr).
\]
Hier bezeichnet $R^{-1}[V]$ eine Menge von Argumenten.
\end{theorem}

\section{Existenz und elementare Konstruktionen}

\begin{theorem}[Funktionsbildung durch eindeutige Beschreibung]
\label{b05guide:thm:description}
Seien $A,B$ Mengen und $\varphi(x,y)$ eine mengentheoretische
Formel, gegebenenfalls mit festen Parametern. Gilt
\[
 \forall x\in A\,\exists!y\in B\,\varphi(x,y),
\]
so existiert genau eine Funktion $f\colon A\to B$ mit
$\varphi(x,f(x))$ für alle $x\in A$. Ihr Graph ist
\[
 f=\{(x,y)\in A\times B\mid\varphi(x,y)\}.
\]
Dies ist ein Aussagenschema für jede solche Formel.
\end{theorem}

\begin{definition}[Konstante Funktion]
Für $b\in B$ ist
\[
 c_{A,b}=A\times\{b\}\colon A\to B,\qquad c_{A,b}(x)=b\quad(x\in A)
\]
die konstante Funktion mit Wert $b$.
Für $A=\varnothing$ ergibt sich unabhängig von $b$ der leere Graph.
\end{definition}

\begin{theorem}[Leere Bereiche und Existenzkriterium]
\label{b05guide:thm:empty}
Für beliebige Mengen $A,B$ gelten:
\begin{enumerate}
 \item Die einzige Funktion $\varnothing\to B$ ist $\varnothing$.
 \item Eine Funktion $A\to\varnothing$ existiert genau dann,
       wenn $A=\varnothing$; dann ist sie der leere Graph.
 \item Eine Funktion $A\to B$ existiert genau dann,
       wenn $A=\varnothing$ oder $B\ne\varnothing$.
 \item Eine Funktion $f\colon A\to B$ hat genau dann einen
       leeren Graphen, wenn $A=\varnothing$.
\end{enumerate}
\end{theorem}

\begin{theorem}[Einelementige Bereiche]
Für $b\in B$ ist $\{(a,b)\}\colon\{a\}\to B$ mit Wert $b$ an $a$.
Jede Funktion auf $\{a\}$ hat diese Gestalt für genau ein $b\in B$.
Für jede Menge $A$ gibt es genau eine Funktion $A\to\{b\}$,
nämlich $c_{A,b}$.
\end{theorem}

\section{Gleichheit und Bestimmtheit}

\begin{theorem}[Definitionsbereich und zulässige Zielmengen]
\label{b05guide:thm:domain}
Aus $f\colon A\to B$ und $f\colon C\to D$ folgt $A=C$.
Für $f\colon A\to B$ und eine beliebige Menge $D$ gilt
\[
 f\colon A\to D
 \quad\Longleftrightarrow\quad
 \forall x\in A\,(f(x)\in D).
\]
Die Zielmenge ist damit durch den Graphen nicht eindeutig festgelegt.
\end{theorem}

\begin{theorem}[Funktionsextensionalität]
\label{b05guide:thm:extensionality}
Für $f\colon A\to B$ und $g\colon C\to D$ gilt
\[
 f=g\quad\Longleftrightarrow\quad
 \bigl(A=C\ \land\ \forall x\in A\,(f(x)=g(x))\bigr).
\]
Die Wertgleichheit auf der rechten Seite wird unter $A=C$ gelesen.
Bei gemeinsamem Definitionsbereich $A$ gilt insbesondere
\[
 f\ne g\quad\Longleftrightarrow\quad
 \exists x\in A\,(f(x)\ne g(x)).
\]
\end{theorem}

\begin{theorem}[Enthaltene Funktionsgraphen]
\label{b05guide:thm:graphinclusion}
Für $f\colon A\to B$ und $g\colon C\to D$ gilt
\[
 f\subseteq g\quad\Longleftrightarrow\quad
 \bigl(A\subseteq C\ \land\ \forall x\in A\,(f(x)=g(x))\bigr).
\]
Bei gleichem Definitionsbereich folgt aus $f\subseteq g$ bereits $f=g$.
\end{theorem}
